\documentclass[]{article}
% Added package start here
\usepackage{array}
\usepackage{geometry} % for margin of paper
\geometry{margin=1.2in, bottom=1.5in}
\usepackage{fancyhdr} % for headers and footers
\usepackage{graphicx} % for image
\graphicspath{ {media/} }
\pagenumbering{gobble} % Delete page number
% End here
\usepackage{lmodern}
\usepackage{amssymb,amsmath}
\usepackage{ifxetex,ifluatex}
\usepackage{fixltx2e} % provides \textsubscript
\ifnum 0\ifxetex 1\fi\ifluatex 1\fi=0 % if pdftex
  \usepackage[T1]{fontenc}
  \usepackage[utf8]{inputenc}
\else % if luatex or xelatex
  \ifxetex
    \usepackage{mathspec}
  \else
    \usepackage{fontspec}
  \fi
  \defaultfontfeatures{Ligatures=TeX,Scale=MatchLowercase}
\fi
% use upquote if available, for straight quotes in verbatim environments
\IfFileExists{upquote.sty}{\usepackage{upquote}}{}
% use microtype if available
\IfFileExists{microtype.sty}{%
\usepackage{microtype}
\UseMicrotypeSet[protrusion]{basicmath} % disable protrusion for tt fonts
}{}
\usepackage{hyperref}
\hypersetup{unicode=true,
            pdfborder={0 0 0},
            breaklinks=true}
\urlstyle{same}  % don't use monospace font for urls
\usepackage{graphicx,grffile}
\makeatletter
\def\maxwidth{\ifdim\Gin@nat@width>\linewidth\linewidth\else\Gin@nat@width\fi}
\def\maxheight{\ifdim\Gin@nat@height>\textheight\textheight\else\Gin@nat@height\fi}
\makeatother
% Scale images if necessary, so that they will not overflow the page
% margins by default, and it is still possible to overwrite the defaults
% using explicit options in \includegraphics[width, height, ...]{}
\setkeys{Gin}{width=\maxwidth,height=\maxheight,keepaspectratio}
\IfFileExists{parskip.sty}{%
\usepackage{parskip}
}{% else
\setlength{\parindent}{0pt}
\setlength{\parskip}{6pt plus 2pt minus 1pt}
}
\setlength{\emergencystretch}{3em}  % prevent overfull lines
\providecommand{\tightlist}{%
  \setlength{\itemsep}{0pt}\setlength{\parskip}{0pt}}
\setcounter{secnumdepth}{0}
% Redefines (sub)paragraphs to behave more like sections
\ifx\paragraph\undefined\else
\let\oldparagraph\paragraph
\renewcommand{\paragraph}[1]{\oldparagraph{#1}\mbox{}}
\fi
\ifx\subparagraph\undefined\else
\let\oldsubparagraph\subparagraph
\renewcommand{\subparagraph}[1]{\oldsubparagraph{#1}\mbox{}}
\fi

\date{}

\begin{document}
% Added information start here
\pagestyle{fancy}
\renewcommand{\headrulewidth}{0pt}
% footer
  \fancyfoot[L]{\footnotesize This material is based upon work supported by the U.S. Department of Energy Office of Science, Advanced Scientific Computing Research and Biological and Environmental Research programs. \begin{flushright} \footnotesize Version 0.2, April 25, 2016 \end{flushright}}
  
\fancypagestyle{empty}{
% header
\fancyhead[C]{\Large {\textbf{What Are Interoperable}}\\ \Large {\textbf{Software Libraries}}\\ \normalsize{Introducing the xSDK}\\ \normalsize {The IDEAS Scientific Software Productivity Project} \\ \small {\href{https://ideas-productivity.org/resources/howtos/}{\emph{ideas-productivity.org/resources/howtos/}}}} 
  \fancyhead[R]{\includegraphics[width = 3 cm, height = 1.8 cm]{ideas_Whatis}}	
}
\thispagestyle{empty}
\textbf{\newline}
\textbf{\newline}
\textbf{\newline}
% Added information ends here.

\textbf{Background:} As extreme-scale computational science increasingly
incorporates multiscale and multiphysics modeling, simulation, and
analysis, the combined use of software developed by independent groups
has become imperative: no single team has resources for the full range
of capabilities needed for predictive science and decision support.
\textbf{Software libraries} have proven effective in providing widely
reusable software that is robust, efficient, and scalable. Moreover,
scientific application codes can employ library design principles to
help manage complexity and achieve good performance, whether the
application software is intended for use in a single context or modest
reuse across applications in the same domain (e.g., as \textbf{domain
components}). While the following discussion uses terminology of
software library interoperability, the concepts also apply to
application-specific domain components.

\textbf{A software library} is a high-quality, encapsulated, documented,
tested, and multiuse software collection that provides functionality
commonly needed by application developers. Key advantages of software
libraries include leverage of library developer expertise and reduced
application coding effort. For example, numerical software libraries
provide easy access to sophisticated mathematical algorithms and
high-performance data structures that have been developed by experts, so
that application users do not need to write this complex code and can
instead focus on their scientific domain software.

Libraries can provide control inversion via abstract interfaces,
call-backs, or similar techniques such that user-defined functionality
can be invoked by the library, for example, a user-defined sparse matrix
multiplication routine. Libraries can also facilitate construction of
related specific objects that provide customizable behavior to improve
performance or flexibility. Moreover, libraries can include
domain-specific software components that are designed to be used by more
than one application.

\textbf{Software library interoperability} refers to the ability of two
or more libraries to be used together in an application code, without
special effort on the part of the user. For simplicity, we discuss
interoperability between two libraries; extension to interoperability
among three or more libraries is conceptually straightforward. Depending
on application needs, various levels of interoperability can be
considered:

\begin{itemize}
\item
  \begin{quote}
  \textbf{Interoperability level 1:} both libraries can be used (side by
  side) in an application
  \end{quote}
\item
  \begin{quote}
  \textbf{Interoperability level 2}: the libraries can exchange data (or
  control data) with each other
  \end{quote}
\end{itemize}

\begin{itemize}
\item
  \begin{quote}
  \textbf{Interoperability level 3}: each library can call the other
  library to perform unique computations
  \end{quote}
\end{itemize}

The simplest case (interoperability level 1) occurs when an application
needs to call two distinct libraries for different functionalities (for
example, an \href{http://www.mpi-forum.org}{\emph{MPI}} library for
message-passing communication and
\href{https://www.hdfgroup.org/HDF5/}{\emph{HDF5}} for data output). As
discussed in
{[}\href{https://figshare.com/articles/Package_Management_Practices_Essential_for_Interoperability_Lessons_Learned_and_Strategies_Developed_for_FASTMath/789055}{\emph{1}},
\href{https://wci.llnl.gov/codes/smartlibs/UCRL-JRNL-208636.pdf}{\emph{2}}{]},
even this basic level of interoperability requires consistency among
libraries to be used within the same application, in terms of compiler,
compiler version/options, and other third-party capabilities. If both
libraries have a dependence on a common third party, the libraries must
be able to use a single common instance of it. For example, more than
one version of the popular SuperLU linear solver library exists, and
interfaces have evolved. If two libraries both use SuperLU, they must be
able to work with the same version of SuperLU. In practice, installing
multiple independently developed packages together can be a tedious
trial-and-error process.

Interoperability level 2 builds on level 1 by enabling conversion, or
encapsulation, and exchange of data between libraries. This level can
simplify use of libraries in sequence by an application. In this case,
the libraries themselves are typically used without internal
modification to support the interoperability.

Interoperability level 3 builds on level 2 by supporting the use of one
library to provide functionality on behalf of another library. This
level of interoperability provides significant value to application
developers because they can access capabilities of additional libraries
through the familiar interfaces of the first library.

%\includegraphics[width=1.29550in,height=1.01354in]{media/image01.png}
\begin{minipage}{0.75\textwidth}

\textbf{The Extreme-Scale Scientific Software Development Kit (xSDK)}
A key aspect of work in the
\href{https://ideas-productivity.org/}{\emph{IDEAS project}} is
development of the \href{http://xsdk.info}{\emph{Extreme-scale
Scientific Software Development Kit (xSDK)}}---a collection of related
and complementary software elements that provide the building blocks,
tools, models, processes, and related artifacts for rapid and efficient
development of high-quality applications.
\end{minipage}
\begin{minipage}{0.3\textwidth}\raggedleft
\includegraphics[width=0.8\textwidth]{xsdk}\raggedleft
\end{minipage}

\textbf{xSDK package compliance standards:} The xSDK addresses
interoperability among the high-performance numerical libraries
\href{https://computation.llnl.gov/project/linear_solvers/}{\emph{hypre}},
\href{https://www.mcs.anl.gov/petsc/}{\emph{PETSc}},
\href{http://crd.lbl.gov/~xiaoye/SuperLU/}{\emph{SuperLU}}, and
\href{https://trilinos.org/}{\emph{Trilinos}}. The xSDK ensures level 1
interoperability for each xSDK library via a full-featured build script
and testing environment and a collection of compliance standards. The
following draft xSDK standards address challenges in interoperability
level 1.

\begin{itemize}
\item
  \begin{quote}
  \textbf{\href{https://docs.google.com/document/d/1DCx2Duijb0COESCuxwEEK1j0BPe2cTIJ-AjtJxt3290/edit\#heading=h.3hcywq83w3x0}{\emph{xSDK
  package compliance standards}}: \emph{A set of minimum requirements}}
  (including topics of configuring, installing, testing, use of MPI,
  portability, contact and version information, open source licensing,
  namespacing, and repository access) that a software package must
  satisfy in order to be considered xSDK compliant. The designation of
  xSDK compliance informs potential users that the package can be easily
  used with other xSDK libraries and components. Also presented are
  \emph{\textbf{recommended standards}} (including topics of public
  repository access, error handling, freeing system resources, and
  library dependencies), which are encouraged but not required.
  \end{quote}
\item
  \begin{quote}
  \textbf{\href{https://docs.google.com/document/d/18028D6nsuhIrCvJnX6c07r8m_Np4SH-aGXMX4svMs1w/edit\#heading=h.3hcywq83w3x0}{\emph{xSDK
  standard configure and CMake options}}: \emph{A standard subset of
  configure and CMake options for xSDK and other HPC packages}} in order
  to make the configuration and installation as efficient as possible on
  standard Linux distributions and Mac OS, as well as on target machines
  at DOE computing facilities
  (\href{http://www.alcf.anl.gov/}{\emph{ALCF}},
  \href{http://www.nersc.gov/}{\emph{NERSC}}, and
  \href{https://www.olcf.ornl.gov/}{\emph{OLCF}}).
  \end{quote}
\end{itemize}

The xSDK collection of software elements commits to compliance with
these guidelines in order to ensure compatibility with other packages
that meet the same standards. The aim is to simplify the combined use of
multiple independently developed software packages and to provide a
foundation for addressing broader issues in interoperability and
performance portability.

\textbf{Deeper levels of xSDK interoperability} involve exchanging,
controlling, and interpreting data, as well as calling routines between
libraries (interoperability levels 2 and 3 described above). Initial
xSDK capabilities of hypre, PETSc, SuperLU, and Trilinos support
interoperability among scalable linear solvers, so that applications can
readily experiment with algorithms across multiple packages, in
combination. Forthcoming companion documents will explain approaches
used for interfaces and adapters between packages as well as work on
interoperability layers for other functionalities. A longer-term goal is
collaboration among members of the HPC community to improve software
interoperability as needed by extreme-scale computational science.\pagebreak 

\textbf{References:}

{[}\href{https://figshare.com/articles/Package_Management_Practices_Essential_for_Interoperability_Lessons_Learned_and_Strategies_Developed_for_FASTMath/789055}{\emph{1}}{]}
Package Management Practices Essential for Interoperability: Lessons
Learned and Strategies Developed for FASTMath, M. C. Miller, L. Diachin,
S. Balay, L. C. McInnes, and B. Smith, First Workshop on Sustainable
Software for Science: Practice and Experiences (WSSSPE), Nov 17,
SC13.

{[}\href{https://wci.llnl.gov/codes/smartlibs/UCRL-JRNL-208636.pdf}{\emph{2}}{]}
Smart Libraries: Best SQE Practices for Libraries with Emphasis on
Scientific Computing, M. C. Miller, J.F. Reus, R.P. Matzke, Q.A. Koziol,
A.P. Cheng, Proceedings of the Nuclear Explosives Code Developer's
Conference, Dec 2004.

This document was prepared by Lois Curfman McInnes, Michael Heroux,
Xiaoye Li, Barry Smith, and Ulrike Yang, with contributions from all
xSDK developers.

\end{document}
